%! AP Statistics — Unit 2, Lesson 2.1 — Guided Notes KEY
\documentclass[10pt]{article}
\usepackage{apstats-article}
\usepackage{apstats-key}

%=============================================================================
\begin{document}

\pageheader{Unit 2, Lesson 2.1}{Guided Notes --- Answer Key}
\begin{tabularx}{\linewidth}{X X X}
Name:\;\ans{Key} & Date:\; & Period:\;
\end{tabularx}

\vspace{0.12in}

\begin{objectivebox}
\small By the end of this lesson, I will be able to\dots\\[4pt]
\ans{identify a statistical question about a possible relationship between two variables;}\\[1pt]
\ans{explain why an apparent association may be random, not real;}\\[1pt]
\ans{explain why association does not imply causation, and give a confounding variable.}
\end{objectivebox}

\vspace{0.1in}

\begin{vocabbox}
\termblanklong{Two-variable data}\hfill\ans{a dataset in which two variables are recorded for each individual}\\
\termblanklong{Association}\hfill\ans{a relationship between two variables --- knowing one gives information about the other}\\
\termblanklong{Apparent association}\hfill\ans{a visible pattern in a sample that may or may not reflect a real relationship}\\
\termblanklong{Causation}\hfill\ans{one variable directly produces a change in another; requires a controlled experiment}\\
\termblanklong{Confounding variable}\hfill\ans{a third variable related to both the explanatory and response variable}\\
\termblanklong{Explanatory variable}\hfill\ans{the variable used to explain or predict the response (on the $x$-axis)}\\
\termblanklong{Response variable}\hfill\ans{the outcome of interest (on the $y$-axis)}
\end{vocabbox}

\vspace{0.1in}

\begin{hookbox}
\small
\begin{enumerate}[leftmargin=1.5em, itemsep=2pt]
  \item \ans{The relationship could be real (students who eat breakfast may be better-rested and
  more focused) or partly random if the sample is small. A large, controlled study would help
  determine which.}

  \item \ans{Age is the lurking variable. Older children both wear larger shoes and have had
  more time to develop reading skills. Shoe size does not cause reading ability.}
\end{enumerate}
\end{hookbox}

\vspace{0.1in}

\begin{notesbox}{From One-Variable to Two-Variable Questions}
\small
\renewcommand{\arraystretch}{1.5}
\begin{tabularx}{\linewidth}{@{} >{\columncolor{sky}\bfseries\color{navy}}p{0.36\linewidth}
                                    X @{}}
\textbf{Unit 1 question (one variable)} & \textbf{Unit 2 question (two variables)} \\
\hline
What is the typical GPA of students at our school?
  & Is there a relationship between sleep and GPA? \\
How many hours of sleep do students get?
  & Do students who play sports sleep more or less than those who don't? \\
\end{tabularx}

\vspace{8pt}
\textbf{The shift:} Unit 2 asks whether two variables \ans{move together (are associated)}
--- whether knowing one variable tells us something about the \ans{other}.
\end{notesbox}

\vspace{0.1in}

\begin{notesbox}{Apparent Associations May Be Random (VAR-1.D.1)}
\small
\begin{multicols}{2}

\textbf{Key idea:}\\
A pattern visible in a \ans{small sample} may disappear when
we collect \ans{more} data.
An apparent association is \textbf{not} guaranteed to
be a real relationship in the \ans{population}.

\vspace{6pt}
\ans{No. With only 10 students the pattern could easily be random.
We need a larger sample and statistical tools to judge whether the association
is real or due to chance.}

\columnbreak

\textbf{What makes an association ``real'' vs.\ ``random''?}
\begin{itemize}[itemsep=3pt]
  \item Larger \ans{sample size} $\to$ less likely to be random
  \item Stronger \ans{association} $\to$ less likely to be random
  \item Consistent across \ans{different} groups $\to$ more credible
\end{itemize}

\end{multicols}
\end{notesbox}

\vspace{0.1in}

\begin{notesbox}{Association vs.\ Causation}
\small
\begin{multicols}{2}

\textbf{Three explanations:}
\begin{enumerate}[leftmargin=1.4em, itemsep=4pt]
  \item $X$ \ans{causes} $Y$
  \item $Y$ \ans{causes} $X$
  \item A third variable \ans{causes} both $X$ and $Y$
        (this is called a \ans{confounding} variable)
\end{enumerate}

\columnbreak

\ans{No. More fire trucks are sent to bigger fires, and bigger fires also cause more damage.
The size of the fire is the confounding variable --- it causes both the number of trucks
dispatched and the amount of property damage.}\\[6pt]
\ans{The size (severity) of the fire.}

\end{multicols}
\end{notesbox}

\vspace{0.1in}

\begin{practicebox}
\small
\vspace{5pt}
\begin{tabularx}{\linewidth}{@{} l X @{}}
\textbf{Two variables:} & \ans{Time on social media (hours/day); life satisfaction score} \\[6pt]
\textbf{Types:}         & \ans{Both quantitative} \\[6pt]
\textbf{Random?}        & \ans{Possibly, especially in a small sample; a large study would help} \\[6pt]
\textbf{Confounder:}    & \ans{Accept any reasonable answer: e.g., mental health status, sleep, academic stress} \\[6pt]
\textbf{Causation?}     & \ans{No --- this is observational data; a randomized experiment would be needed} \\[6pt]
\end{tabularx}
\end{practicebox}

\vspace{0.1in}

\begin{teachernote}
\textbf{Common issues in notes:}
\begin{itemize}[itemsep=3pt]
  \item Students who say ``the sample is too small'' for the coffee-GPA example are on the right track
        --- push them to articulate \emph{why}: with $n=10$ the pattern is very likely random.
  \item For the fire-trucks example, the most common wrong answer is ``the association is negative''
        (students thinking fewer trucks = less damage). Redirect: the \emph{direction} of the observed
        association is positive; the issue is what causes it.
  \item The social media practice item is intentionally open-ended. Accept any plausible confounding
        variable; use class discussion to highlight that the best confounders are related to \emph{both} variables.
\end{itemize}
\end{teachernote}

\end{document}
